\documentclass[12pt,a4paper]{article}

\usepackage[margin=1in]{geometry}
\usepackage[T1]{fontenc}
\usepackage[utf8]{inputenc}
\usepackage{lmodern}
\usepackage{setspace}
\usepackage{booktabs}
\usepackage{longtable}
\usepackage{array}
\usepackage{hyperref}
\usepackage{xcolor}
\usepackage{enumitem}
\usepackage{amsmath}
\usepackage{fancyhdr}

\hypersetup{
  colorlinks=true,
  linkcolor=blue,
  urlcolor=blue,
  pdftitle={AI Emergency Command System - Technical Report},
  pdfauthor={Solution Challenge Team}
}

\setstretch{1.15}
\pagestyle{fancy}
\fancyhf{}
\lhead{AI Emergency Command System}
\rhead{Technical Report}
\cfoot{\thepage}

\title{\textbf{AI Emergency Command System}\\\large End-to-End Technical Report}
\author{Solution Challenge Project Team}
\date{April 25, 2026}

\begin{document}
\maketitle
\tableofcontents
\newpage

\section{Executive Summary}
This project implements an AI-assisted emergency operations pipeline that converts unstructured incident reports into structured operational intelligence. The system combines:
\begin{itemize}[leftmargin=1.5em]
  \item a FastAPI backend,
  \item a Random Forest machine learning model for escalation prediction,
  \item a Gemini large language model for triage and tactical generation tasks,
  \item and a Next.js dashboard with Google Maps visualization.
\end{itemize}

The workflow supports four primary actions:
\begin{enumerate}[leftmargin=1.5em]
  \item analyze raw emergency text,
  \item predict escalation risk,
  \item generate auto-dispatch recommendations,
  \item generate a tactical briefing for field command.
\end{enumerate}

\section{Problem Statement and Objectives}
Emergency command workflows often receive chaotic, incomplete, and time-sensitive text input. The project objective is to reduce decision latency by building a decision-support stack that:
\begin{itemize}[leftmargin=1.5em]
  \item extracts structured incident fields from natural-language reports,
  \item predicts whether the threat is likely to escalate,
  \item recommends initial dispatch resource allocation,
  \item and provides a concise command briefing.
\end{itemize}

\section{System Architecture}
\subsection{High-Level Architecture}
\begin{verbatim}
Frontend (Next.js + React + Google Maps)
    |
    +--> POST /analyze-input       (LLM: triage + extraction)
    +--> POST /predict-escalation  (ML: RandomForestClassifier)
    +--> POST /auto-dispatch       (LLM: dispatch planning)
    +--> POST /generate-briefing   (LLM: tactical summary)
    |
Backend (FastAPI + scikit-learn + Gemini SDK)
\end{verbatim}

\subsection{Backend Components}
\begin{itemize}[leftmargin=1.5em]
  \item \textbf{Framework:} FastAPI with async endpoints.
  \item \textbf{Startup lifecycle:} ML model is trained once during app startup.
  \item \textbf{CORS:} enabled for all origins to support local frontend integration.
  \item \textbf{Data source:} CSV dataset at \texttt{backend/data/escalation\_dataset.csv}.
\end{itemize}

\subsection{Frontend Components}
\begin{itemize}[leftmargin=1.5em]
  \item \textbf{Input Panel:} receives emergency narrative text.
  \item \textbf{Intel Panel:} displays triage category, risk score, dispatch recommendation, and final briefing.
  \item \textbf{Map View:} plots incidents using Google Maps markers with severity-based color coding.
  \item \textbf{Status Bar:} shows pipeline stage and system status.
\end{itemize}

\section{Detailed Processing Flow}
\subsection{Flow Logic}
\begin{enumerate}[leftmargin=1.5em]
  \item User submits raw emergency text from dashboard.
  \item Backend endpoint \texttt{/analyze-input} prompts Gemini to return strict JSON fields:
  incident type, location, severity level, priority score, triage category, and reason.
  \item Frontend maps extracted data to ML feature schema and calls \texttt{/predict-escalation}.
  \item Backend model returns threat growth probability (\%) and risk-band interpretation.
  \item Frontend calls \texttt{/auto-dispatch} to obtain tactical resource orders.
  \item Frontend calls \texttt{/generate-briefing} to get command-ready text summary.
  \item Results are rendered in the intelligence panel and mapped as an active incident marker.
\end{enumerate}

\subsection{Operational Pipeline Table}
\begin{longtable}{p{0.23\linewidth} p{0.22\linewidth} p{0.24\linewidth} p{0.23\linewidth}}
\toprule
\textbf{Stage} & \textbf{Input} & \textbf{Engine} & \textbf{Output} \\
\midrule
Analyze Input & Raw incident text & Gemini prompt with START triage constraints & Structured incident JSON \\
Predict Escalation & Structured features (7 fields) & RandomForestClassifier & Threat growth probability + risk statement \\
Auto Dispatch & Incident JSON + resource list & Gemini dispatch prompt & Actionable dispatch order \\
Generate Briefing & Incident + prediction + dispatch & Gemini briefing prompt & Tactical command briefing \\
\bottomrule
\end{longtable}

\section{ML Model Design and Training}
\subsection{Model Type}
\begin{itemize}[leftmargin=1.5em]
  \item \textbf{Algorithm:} \texttt{RandomForestClassifier}
  \item \textbf{Library:} scikit-learn
  \item \textbf{Hyperparameters in implementation:}
  \begin{itemize}[leftmargin=1.5em]
    \item \texttt{n\_estimators = 100}
    \item \texttt{random\_state = 42}
  \end{itemize}
\end{itemize}

\subsection{Feature Schema}
The model is trained on 7 categorical features:
\begin{enumerate}[leftmargin=1.5em]
  \item incident\_type
  \item location\_type
  \item time\_of\_day
  \item severity\_level
  \item people\_trapped
  \item hazardous\_material
  \item resource\_availability
\end{enumerate}
Target label: \texttt{escalation} (yes/no).

\subsection{Preprocessing}
Each categorical feature and the target variable are encoded with \texttt{LabelEncoder}. The backend also includes unknown-category fallback logic during inference (unseen values map to encoded index 0).

\subsection{Dataset Profile}
Measured from the current dataset:
\begin{itemize}[leftmargin=1.5em]
  \item Total records: 500
  \item Target distribution: 273 yes, 227 no
\end{itemize}

\section{Measured ML Performance}
The following metrics were computed from the current dataset using a reproducible evaluation script.

\subsection{Evaluation Setup}
\begin{itemize}[leftmargin=1.5em]
  \item Train-test split: 80:20, stratified, \texttt{random\_state = 42}
  \item Additional validation: 5-fold stratified cross-validation
\end{itemize}

\subsection{Results}
\begin{table}[h!]
\centering
\begin{tabular}{lr}
\toprule
\textbf{Metric} & \textbf{Value} \\
\midrule
Training Accuracy & 0.9500 \\
Test Accuracy & 0.6600 \\
Precision (positive class) & 0.6981 \\
Recall (positive class) & 0.6727 \\
F1 Score (positive class) & 0.6852 \\
5-Fold CV Mean Accuracy & 0.7020 \\
5-Fold CV Std Dev & 0.0631 \\
\bottomrule
\end{tabular}
\caption{Escalation model performance metrics}
\end{table}

\subsection{Confusion Matrix (Test Split)}
\[
\begin{bmatrix}
29 & 16 \\
18 & 37
\end{bmatrix}
\]
where rows are true labels [no, yes] and columns are predicted labels [no, yes].

\subsection{Interpretation}
\begin{itemize}[leftmargin=1.5em]
  \item The high training accuracy (0.95) compared to test accuracy (0.66) suggests overfitting risk.
  \item Cross-validation mean around 0.70 indicates moderate generalization.
  \item Current model is useful for risk signaling, but should not be treated as a sole decision authority.
\end{itemize}

\section{LLM Integration Details}
\subsection{Model Used in Backend}
\begin{itemize}[leftmargin=1.5em]
  \item Gemini model ID configured in backend code: \texttt{gemini-2.5-flash}
  \item SDK: \texttt{google-generativeai}
\end{itemize}

\subsection{LLM Responsibilities}
\begin{enumerate}[leftmargin=1.5em]
  \item \textbf{Triage analysis:} structured extraction and START category assignment.
  \item \textbf{Dispatch planning:} resource allocation recommendation from incident context.
  \item \textbf{Briefing generation:} concise command-ready tactical summary.
\end{enumerate}

\subsection{Prompt Engineering Strategy}
\begin{itemize}[leftmargin=1.5em]
  \item Hard output schema for JSON tasks to reduce parser failures.
  \item Domain constraints embedded in prompt (hazmat, trapped persons, severity escalation rules).
  \item Post-processing helper extracts valid JSON from either plain text or fenced responses.
\end{itemize}

\subsection{Implementation Note}
The UI status text currently references ``GEMINI 2.0 FLASH'', while backend runtime configuration uses \texttt{gemini-2.5-flash}. This should be synchronized for consistency.

\section{Google Maps API Integration}
\subsection{Library and Key Usage}
\begin{itemize}[leftmargin=1.5em]
  \item React integration library: \texttt{@vis.gl/react-google-maps}
  \item API key source: \texttt{NEXT\_PUBLIC\_GOOGLE\_MAPS\_API\_KEY}
  \item Map components used: \texttt{APIProvider}, \texttt{Map}, \texttt{AdvancedMarker}, \texttt{Pin}
\end{itemize}

\subsection{Map Behavior}
\begin{itemize}[leftmargin=1.5em]
  \item Default map center is configured to Chennai, India.
  \item Incidents are visualized as markers with severity-color mapping:
  low (green), medium (amber), high (red), critical (violet).
  \item Marker jitter/offset is applied to reduce overlap when multiple incidents are plotted near the same area.
  \item If key is not configured, the UI gracefully displays a setup instruction panel.
\end{itemize}

\section{API Endpoints and Contracts}
\begin{longtable}{p{0.22\linewidth} p{0.73\linewidth}}
\toprule
\textbf{Endpoint} & \textbf{Purpose} \\
\midrule
\texttt{POST /analyze-input} & Converts raw report text into structured fields and START triage classification. \\
\texttt{POST /predict-escalation} & Predicts threat growth percentage and risk message using the trained Random Forest model. \\
\texttt{POST /auto-dispatch} & Generates resource deployment action and tactical reason via Gemini. \\
\texttt{POST /generate-briefing} & Produces concise tactical briefing text for commanders. \\
\texttt{GET /health} & Service health and model load status check. \\
\bottomrule
\end{longtable}

\section{How It Was Built (Implementation Narrative)}
\begin{enumerate}[leftmargin=1.5em]
  \item Defined operational target: transform narrative emergency text into actionable intelligence.
  \item Built backend API skeleton with FastAPI and clear endpoint separation by function.
  \item Added Gemini integration for language-heavy tasks requiring contextual reasoning.
  \item Designed and trained a tabular ML classifier on escalation dataset for quantitative risk scoring.
  \item Implemented startup training lifecycle so model is loaded before first prediction request.
  \item Built frontend dashboard with sequential orchestration across all four backend stages.
  \item Added map visualization for command situational awareness.
  \item Added status and reasoning panels for transparency and operator trust.
\end{enumerate}

\section{Current Limitations and Risks}
\begin{itemize}[leftmargin=1.5em]
  \item Categorical fallback for unseen labels currently defaults to class index 0; this may bias outputs.
  \item Training and inference are based on synthetic/limited tabular features, reducing real-world nuance.
  \item No explicit model versioning, drift detection, or calibration pipeline yet.
  \item LLM output quality depends on prompt reliability and API availability.
\end{itemize}

\section{Recommended Next Improvements}
\begin{enumerate}[leftmargin=1.5em]
  \item Add held-out validation set tracking and periodic retraining with versioned artifacts.
  \item Add probability calibration and threshold tuning for mission-dependent false-negative control.
  \item Replace random map jitter with geocoded incident coordinates for true geospatial accuracy.
  \item Add request logging and trace IDs to improve observability and post-incident auditability.
  \item Align displayed LLM version text with backend runtime configuration.
\end{enumerate}

\section{Conclusion}
The implemented system already demonstrates a full operational loop from incident intake to command briefing, combining both deterministic ML scoring and LLM-assisted reasoning. The measured accuracy indicates moderate predictive capability, suitable for decision support. With additional data, calibration, and observability hardening, this architecture can evolve into a robust emergency command intelligence platform.

\end{document}
